
\newcommand{\DataSummaryTable}{%
\begin{table}[!htbp]
\centering
\caption{Prepared ChEMBL 36 activity dataset and fold-0 split.}
\label{tab:data-summary}
\begin{tabular}{lr}
\toprule
Quantity & Value \\
\midrule
Raw activity rows queried & 3,501,718 \\
Canonical molecules after filtering and deduplication & 1,369,515 \\
Binary activity threshold & pChEMBL $\geq$ 7.0 \\
Active molecules & 520,188 \\
Temporal holdout molecules & 274,541 \\
Fold-0 training molecules & 875,979 \\
Fold-0 internal-validation molecules & 218,995 \\
\bottomrule
\end{tabular}
\end{table}
}

\newcommand{\ModelPerformanceTable}{%
\begin{table}[!htbp]
\centering
\caption{Activity-model performance on ChEMBL 36 holdouts.}
\label{tab:model-performance}
\begin{tabular}{lrrrr}
\toprule
Split & RMSE & $R^2$ & Spearman & EF@1\% \\
\midrule
Internal validation & 0.8466 & 0.5992 & 0.7674 & 2.7331 \\
Temporal holdout & 1.1615 & 0.2473 & 0.5171 & 2.4359 \\
\bottomrule
\end{tabular}
\end{table}
}

\newcommand{\BaselineSanityTable}{%
\begin{table}[!htbp]
\centering
\caption{Selected baseline sanity checks, not a full benchmark suite. Random forest and Chemprop use sampled ChEMBL 36 splits for context only; they are not full-dataset, same-protocol baselines. The current model and previous so\_f4 comparison use the full future prediction artifacts.}
\label{tab:baselines}
\begin{tabular}{llrrrr}
\toprule
Model & Evaluation rows & RMSE & $R^2$ & Spearman & EF@1\% \\
\midrule
Current 2D MLP & Future full (274,541) & 1.1615 & 0.2473 & 0.5171 & 2.4359 \\
Previous so\_f4 & Future full (274,541) & 1.1914 & 0.2080 & 0.4806 & 2.1008 \\
Random forest Morgan & Future sample (10,000) & 1.2236 & 0.1736 & 0.4203 & 1.9859 \\
Chemprop D-MPNN & Future sample (10,000) & 1.2164 & 0.1834 & 0.4337 & 1.8100 \\
\bottomrule
\end{tabular}
\end{table}
}

\newcommand{\BaceOverlapTable}{%
\begin{table}[!htbp]
\centering
\caption{BACE external-source performance before and after fold-0 training-overlap control. Primary interpretation should use the strict scaffold-disjoint row.}
\label{tab:bace-overlap}
\begin{tabular}{lrrrrr}
\toprule
Subset & $n$ & ROC AUC & RMSE & Spearman & EF@1\% \\
\midrule
Full, provenance only & 1513 & 0.8304 & 0.9330 & 0.6958 & 2.0527 \\
Structure-disjoint & 1192 & 0.8035 & 1.0052 & 0.6641 & 2.1286 \\
Scaffold-disjoint & 962 & 0.7626 & 1.0370 & 0.6047 & 2.0253 \\
\bottomrule
\end{tabular}
\end{table}
}

\newcommand{\EgfrOverlapTable}{%
\begin{table}[!htbp]
\centering
\caption{EGFR/CHEMBL203 label-hidden operational sensitivity after fold-0 training-overlap control. This same-source replay is not independent external validation or a natural prospective library.}
\label{tab:egfr-overlap}
\begin{tabular}{lrrrrr}
\toprule
Subset & $n$ & ROC AUC & Spearman & EF@1\% & Hit@10\% \\
\midrule
Full, workflow provenance & 500 & 0.9582 & 0.7001 & 2.0000 & 1.0000 \\
Structure-disjoint & 330 & 0.9754 & 0.7479 & 2.0000 & 1.0000 \\
Scaffold-disjoint & 223 & 0.9769 & 0.7773 & 2.1038 & 1.0000 \\
\bottomrule
\end{tabular}
\end{table}
}

\newcommand{\DecisionReplayTable}{%
\begin{table}[!htbp]
\centering
\caption{Strict BACE decision-layer replay on a 100-molecule virtual batch. The risk-aware order changes the front of the list rather than uniformly dominating raw activity sorting.}
\label{tab:decision-replay}
\small
\begin{tabular}{lrrrr}
\toprule
Order & Hit@10 & Mean pIC50@10 & Hit@20 & Mean pIC50@20 \\
\midrule
Raw activity & 1.0000 & 8.2303 & 0.9000 & 7.6705 \\
Decision (triage $\rightarrow$ priority score) & 0.9000 & 7.9546 & 0.8500 & 7.6392 \\
\bottomrule
\end{tabular}
\end{table}
}

\newcommand{\OperationalABTable}{%
\begin{table}[!htbp]
\centering
\caption{Retrospective operational A/B controls on the BACE scaffold-disjoint subset. Random and scaffold-diversity rows are averaged over five seeds.}
\label{tab:operational-ab}
\begin{tabular}{llrr}
\toprule
Top N & Selection rule & Hit rate & Mean pIC50 \\
\midrule
10 & Ranking score & 1.0000 & 8.3516 \\
10 & Random & 0.4400 $\pm$ 0.1342 & 6.6667 $\pm$ 0.3905 \\
10 & Scaffold diversity & 0.4200 $\pm$ 0.1304 & 6.2401 $\pm$ 0.2620 \\
50 & Ranking score & 1.0000 & 8.2026 \\
50 & Random & 0.5240 $\pm$ 0.0654 & 6.8378 $\pm$ 0.1919 \\
50 & Scaffold diversity & 0.4400 $\pm$ 0.0678 & 6.4208 $\pm$ 0.1913 \\
100 & Ranking score & 0.8900 & 7.8911 \\
100 & Random & 0.4740 $\pm$ 0.0541 & 6.6163 $\pm$ 0.1467 \\
100 & Scaffold diversity & 0.4440 $\pm$ 0.0483 & 6.4902 $\pm$ 0.1531 \\
\bottomrule
\end{tabular}
\end{table}
}

\newcommand{\ErrorCasesTable}{%
\begin{table}[!htbp]
\centering
\caption{High-ranked inactive cases from the BACE scaffold-disjoint subset. Hashes are reported instead of full structures in the manuscript; full public benchmark rows remain in the reproducibility artifact.}
\label{tab:error-cases}
\small
\begin{tabular}{lrrrr}
\toprule
Case & Rank & True pIC50 & Predicted pIC50 & Strict-subset scaffold frequency \\
\midrule
BACE\_SD\_052 & 52 & 6.8239 & 7.7409 & 3 \\
BACE\_SD\_057 & 57 & 5.4401 & 7.7257 & 1 \\
BACE\_SD\_060 & 60 & 6.7595 & 7.7065 & 41 \\
\bottomrule
\end{tabular}
\end{table}
}

\newcommand{\StrategyDifferenceTable}{%
\begin{table}[!htbp]
\centering
\caption{Strategy-difference summary on the BACE scaffold-disjoint subset and the strict 100-molecule BACE replay. Top-10 overlap compares frozen selection rules on the same scaffold-disjoint subset; the bucket panel summarizes the label-hidden decision replay and is not a prospective risk-distribution claim.}
\label{tab:strategy-difference}
\footnotesize
\setlength{\tabcolsep}{3pt}
\begin{tabular}{>{\raggedright\arraybackslash}p{0.38\linewidth}p{0.19\linewidth}rr}
\toprule
Item & Context/share & Overlap/hit & Jaccard/pIC50 \\
\midrule
\multicolumn{4}{l}{\textit{Selection-rule contrast}} \\
\midrule
Activity rank vs. Scaffold diversity & BACE SD & 0 & 0.0000 \\
Activity rank vs. Random seed 2026 & BACE SD & 0 & 0.0000 \\
Random seed 2026 vs. Scaffold diversity & BACE SD & 0 & 0.0000 \\
\midrule
\multicolumn{4}{l}{\textit{Decision buckets in strict 100-molecule replay}} \\
\midrule
Priority & 0.1100 & 0.9091 & 8.0637 \\
Watch & 0.1800 & 0.8333 & 7.3748 \\
Low & 0.7100 & 0.4225 & 6.1861 \\
Review & 0.0000 & -- & -- \\
\bottomrule
\end{tabular}
\end{table}
}

\newcommand{\DecisionMatrixTable}{%
\begin{table}[!htbp]
\centering
\caption{Decision matrix used for review guidance. The matrix is a reporting abstraction, not an automatic experimental decision.}
\label{tab:decision-matrix}
\begin{tabular}{lll}
\toprule
Activity evidence & Domain/uncertainty evidence & Suggested review action \\
\midrule
High & Narrow interval or in-domain & First review priority \\
High & Wide interval or out-of-domain & Risk review before advancement \\
Moderate/low & Narrow interval or in-domain & Reserve or diversity control \\
Moderate/low & Wide interval or out-of-domain & Defer without external support \\
\bottomrule
\end{tabular}
\end{table}
}

\newcommand{\CostFramingTable}{%
\begin{table}[!htbp]
\centering
\caption{Assay-budget framing example for library compression. The values illustrate accounting under a fixed Top-10\% budget and are not a demonstrated project-specific savings claim.}
\label{tab:cost-framing}
\begin{tabular}{lr}
\toprule
Quantity & Value \\
\midrule
Initial library size & 10,000 \\
Compression fraction & 10\% \\
First-round tests after compression & 1,000 \\
Initial tests avoided & 9,000 \\
BACE scaffold-disjoint active recall at Top 10\% & 0.1789 \\
\bottomrule
\end{tabular}
\end{table}
}

\newcommand{\VirtualBatchTable}{%
\begin{table}[!htbp]
\centering
\caption{Virtual-batch ranking at batch size 1,000 across five library-shuffle seeds. Variability reflects batch composition, not model retraining.}
\label{tab:batch-ranking}
\small
\setlength{\tabcolsep}{4pt}
\begin{tabular}{lrrrr}
\toprule
Split & EF@1\% & EF@5\% & EF@10\% & NDCG@10 \\
\midrule
Internal validation & 2.7326 $\pm$ 0.0027 & 2.6789 $\pm$ 0.0013 & 2.6108 $\pm$ 0.0007 & 0.9510 $\pm$ 0.0000 \\
Temporal holdout & 2.4398 $\pm$ 0.0080 & 2.1513 $\pm$ 0.0035 & 1.9534 $\pm$ 0.0012 & 0.8811 $\pm$ 0.0002 \\
BACE scaffold-disjoint & 2.0253 $\pm$ 0.0000 & 2.0253 $\pm$ 0.0000 & 1.7956 $\pm$ 0.0000 & 0.9143 $\pm$ 0.0000 \\
\bottomrule
\end{tabular}
\vspace{2pt}
\begin{minipage}{0.92\linewidth}\footnotesize For BACE scaffold-disjoint ($n=962$), the batch-size 1,000 row corresponds to one near-complete virtual batch under each shuffle seed; the zero standard deviation is therefore a consequence of discrete binning rather than a formatting artifact.\end{minipage}
\end{table}
}

\newcommand{\BootstrapTable}{%
\begin{table}[!htbp]
\centering
\caption{Paired bootstrap 95\% intervals over frozen prediction rows. These intervals quantify fixed-prediction uncertainty and exclude retraining variability.}
\label{tab:bootstrap}
\begin{tabular}{lrrrr}
\toprule
Dataset & Rows & RMSE & Spearman & ROC AUC \\
\midrule
Internal validation & 218,995 & 0.8432--0.8497 & 0.7653--0.7696 & 0.8838--0.8867 \\
Temporal holdout & 274,541 & 1.1585--1.1648 & 0.5143--0.5197 & 0.7513--0.7550 \\
BACE scaffold-disjoint & 962 & 0.9861--1.0895 & 0.5598--0.6452 & 0.7315--0.7906 \\
EGFR scaffold-disjoint & 223 & 2.1179--2.3440 & 0.7306--0.8123 & 0.9596--0.9899 \\
\bottomrule
\end{tabular}
\end{table}
}

\newcommand{\ConformalTable}{%
\begin{table}[!htbp]
\centering
\caption{Split-conformal interval coverage.}
\label{tab:conformal}
\begin{tabular}{lrrr}
\toprule
Split & 90\% coverage & 80\% coverage & 90\% width \\
\midrule
Calibration & 0.9000 & 0.8000 & 2.7686 \\
Temporal holdout & 0.7799 & 0.6423 & 2.7686 \\
\bottomrule
\end{tabular}
\end{table}
}
